% Part: incompleteness
% Chapter: incompleteness-provability
% Section: provability-conditions

\documentclass[../../../include/open-logic-section]{subfiles}

\begin{document}

\olfileid[hi]{inc}{inp}{prc}

\olsection{$\Th{PA}$ के लिए \usetoken{S}{derivability} प्रतिबंध}

पियानो अंकगणित, अथवा $\Th{PA}$, वह सिद्धांत है जो सभी !!{formula}s के लिए
आगमन-अभिगृहीत जोड़कर $\Th{Q}$ का विस्तार करता है। दूसरे शब्दों में,
प्रत्येक !!{formula}~$!A$ के लिए $\Th{Q}$ में इस रूप का अभिगृहीत जोड़ते हैं:
\[
(!A(0) \land \lforall[x][(!A(x) \lif !A(x'))]) \lif \lforall[x][!A(x)]
\]
ध्यान दें कि यह वास्तव में एक \emph{योजना} है, अर्थात् अनंततः अनेक
अभिगृहीत (और पता चलता है कि $\Th{PA}$ परिमिततः अभिगृहीतीकरणीय {\em नहीं}
है)। किंतु क्योंकि प्रभावी ढंग से निर्धारित किया जा सकता है कि चिह्नों की
कोई शृंखला आगमन-अभिगृहीत का उदाहरण है या नहीं, इसलिए $\Th{PA}$ के
अभिगृहीतों का समुच्चय संगणनीय है। $\Th{PA}$,~$\Th{Q}$ से कहीं अधिक समर्थ
सिद्धांत है। उदाहरणतः सामान्य ढंग से आगमन का उपयोग करके सरलता से सिद्ध
किया जा सकता है कि जोड़ और गुणन क्रमविनिमेय हैं। वास्तव में अधिकांश
परिमितीय संख्या-सैद्धांतिक और संयोजनात्मक तर्क~$\Th{PA}$ में किए जा सकते
हैं।

चूँकि $\Th{PA}$ संगणनीयतः अभिगृहीतीकृत है, !!{derivability} विधेय
$\Prf[\Th{PA}](x,y)$ संगणनीय है और इसलिए~$\Th{Q}$ में (और इस प्रकार
$\Th{PA}$ में) निरूपित होता है। पहले की तरह संबंध का निरूपण करने वाले
सूत्र को हम $\OPrf[\Th{PA}](x,y)$ लिखेंगे। $\OProv[\Th{PA}](y)$ को सूत्र
$\lexists[x][\Prf[\Th{PA}](x,y)]$ मानें, जो सहजतः कहता है, ``$y$,
$\Th{PA}$ के अभिगृहीतों से !!{derivable} है।'' हमें $\Th{Q}$ के अभिगृहीतों
से कुछ अधिक की आवश्यकता इसलिए है कि प्रयुक्त सिद्धांत इस
!!{derivability} विधेय के बारे में कुछ मूलभूत तथ्यों को !!{derive} करने के
लिए पर्याप्त प्रबल होना चाहिए। वास्तव में हमें निम्न तथ्य चाहिए:
\begin{enumerate}
\item[P1.] यदि $\Th{PA} \Proves !A$, तो
  $\Th{PA} \Proves \OProv[\Th{PA}](\gn{!A})$।
\item[P2.] सभी !!{formula}s $!A$ और $!B$ के लिए,
  \[
  \Th{PA} \Proves \OProv[\Th{PA}](\gn{!A \lif !B}) \lif
  (\OProv[\Th{PA}](\gn{!A}) \lif \OProv[\Th{PA}](\gn{!B})).
  \]
\item[P3.] प्रत्येक !!{formula}~$!A$ के लिए,
  \[
  \Th{PA} \Proves \OProv[\Th{PA}](\gn{!A})
  \lif \OProv[\Th{PA}](\gn{\OProv[\Th{PA}](\gn{!A})}).
  \]
\end{enumerate}
इन तीन गुणों की पुष्टि करने का एकमात्र ढंग !!{formula}
$\OProv[\Th{PA}](y)$ का सावधानी से वर्णन करना और संबद्ध औपचारिक
!!{derivation}s का वर्णन करने के लिए $\Th{PA}$ के अभिगृहीतों का उपयोग करना
है। प्रतिबंध (1) और~(2) सरल हैं; वास्तव में प्रतिबंध~(3) में काम लगता है।
(सोचें कि इसमें किस प्रकार का काम लगेगा \dots।) विवरण पूरा करना श्रमसाध्य
और अरुचिकर होगा, इसलिए यहाँ हम आपसे विश्वास करने को कहेंगे कि $\Th{PA}$
में ऊपर सूचीबद्ध तीन गुण हैं। $\OProv[\Th{PA}](y)$ का कोई उचित चुनाव यह भी
संतुष्ट करेगा:
\begin{enumerate}
\item[P4.] यदि $\Th{PA} \Proves \OProv[\Th{PA}](\gn{!A})$, तो
  $\Th{PA} \Proves !A$।
\end{enumerate}
किंतु हमें इस तथ्य की आवश्यकता नहीं होगी।

\begin{digress}
संयोग से, ग्योडेल (G\"odel) भी उसी ढंग से आलसी थे जैसे हम अभी हैं। 1931 के शोधपत्र
के अंत में वे द्वितीय अपूर्णता-प्रमेय के प्रमाण की रूपरेखा देते हैं और बाद
के शोधपत्र में विवरण देने का वचन देते हैं। उन्होंने उसे कभी पूरा नहीं किया;
क्योंकि तर्क समझने वाले सभी लोगों को विश्वास था कि उसे पूरा किया जा सकता
है, उन्हें विवरण भरने की आवश्यकता नहीं पड़ी।
\end{digress}

\end{document}
